\noindent A translation judge can prefer one candidate to another and still accept a serious error. We therefore examined Chinese news sentences with professional judgments collected independently of this study. For each retained sentence, at least three raters marked one English translation error-free and every rater marked a major accuracy error in another. The original raters could consult the document. Our models saw the sentence alone. The 103 qualifying sentences came from 57 documents in the WMT 2020 news test.\cite{freitagmqm}

\displayfigure{Expert_news_translation.pdf}{Free visual writing and the reliability of its judges}{a, Original source and answer pixels for the first news item in fixed order whose two English translations fit the current writer. The writer repeats its phrase about years. b, Cases where each assessor accepted the error-free translation and rejected the major error. The comparative condition required that decision in both presented orders. Each of the two translations was seen once in isolation in the separate condition. The professional annotations apply to the submitted WMT translations. The displayed ILM answer was generated later.}

\begin{multicols}{2}
\section*{What the current writer can read}
The writer matches each input cell against a fixed library of printed shapes. Of 103 sources, 48 contained a shape outside the admitted Chinese inventory and one exceeded its length limit. We passed all 54 representable sources as pixels. The writer stopped on 52; eight isolated replays reproduced the same pixels and stopping decisions. Only nine cases admitted both complete English alternatives. It preferred the professionally accepted translation on five. Exchanging the Chinese source and recomputing the draft left all nine choices unchanged.
Both assessors rejected the complete meaning of every stopped independent answer in separate single-candidate grading. In the comparative condition each judge reproduced the professional choice for 67/103 segments. When grading the two translations independently, the counts fell to 24 and 40. The single-candidate judges accepted 96 and 94 of 103 professionally approved translations, yet rejected just 28 and 46 of the 103 translations containing a major accuracy error. A positive single-candidate vote is therefore weak evidence of faithfulness on this panel.
\section*{Independent verification of natural source edits}
From one million original OPUS-100 training pairs, 7,081 disjoint families shared a one-word source contrast and 3,413 met four bidirectional likelihood margins. A first model individually confirmed 818 of 2,048 families. The second confirmed all eight decisions on 88 of 128 reserved families, below the registered minimum of 112. Thirty-four of 40 unconfirmed families had an order-sensitive vote. The proposed matched training comparison received zero updates. \repo{references/opus_paired_supervision_result.md}{Registered training decision}. \repo{references/human_translation_calibration_result.md}{Individual news judgments and visual outputs}.
\section*{Reading pixels beyond the original repertoire}
The existing visual network can pool an unfamiliar shape directly. With phrase hints restricted to contiguous recognised runs, it admitted another 48 sources. Hiding the new ink while preserving cell positions changed 47 answers. Both versions received zero complete-meaning credits. A separate pretrained Qwen3-VL-8B read all 103 source images and returned text. The single-candidate grader credited 97 responses from pixels and 102 from text. Its poor rejection rate for major errors limits the interpretation of the absolute scores.

\end{multicols}
